%% ma: L12-B14-0103
%% lop: 12
%% chuong: 5
%% bai: 14
%% dang: 12.14.30
%% loai: TLN
%% muc: VDC
%% dap_an: "27"
%% nguon: "Ảnh giáo viên gửi 07/11/2026 (Tư Duy Mở - tuduymo.com, đề tự luyện Phương trình mặt phẳng, đề lập phương trình mặt phẳng), Câu 21"
%% kien_thuc: "Bài 14 (phương trình mặt phẳng theo đoạn chắn), Bài 2 (giá trị lớn nhất của hàm số bằng đạo hàm), hàm số bậc hai (lớp 10)"
%% trang_thai: chinh-thuc
%% ghi_chu: "Tài liệu gốc không kèm đáp án; đáp án do Claude giải và kiểm bằng sympy. Lời giải phổ biến dùng bất đẳng thức Cô-si ba số (ngoài SGK KNTT); lời giải ở đây chỉ dùng hàm bậc hai và đạo hàm. Dạng đã duyệt 07/11/2026."
\begin{ex}
Trong không gian với hệ trục tọa độ $Oxyz$, gọi mặt phẳng $(P)\colon ax+by+cz+3=0$ là mặt phẳng cắt ba trục tọa độ $Ox$, $Oy$, $Oz$ lần lượt tại ba điểm $A$, $B$, $C$. Trong đó $a$, $b$, $c$ là ba số thực dương thỏa mãn hệ thức $a+2b+3c=3$. Hãy xác định giá trị nhỏ nhất của thể tích tứ diện $OABC$.
\shortans{27}
\loigiai{$A\left(-\dfrac3a;0;0\right)$, $B\left(0;-\dfrac3b;0\right)$, $C\left(0;0;-\dfrac3c\right)$, $V=\dfrac16\cdot\dfrac{27}{abc}=\dfrac{9}{2abc}$.
Với $c$ cố định ($0<c<1$): $a+2b=3-3c$, $ab=\dfrac12\cdot a\cdot2b\le\dfrac{(3-3c)^2}{8}$ (dấu bằng khi $a=2b$).
Xét $f(c)=\dfrac{c(3-3c)^2}{8}$ trên $(0;1)$: $f'(c)=\dfrac{9(1-c)(1-3c)}{8}$, $f$ lớn nhất tại $c=\dfrac13$, $f\left(\dfrac13\right)=\dfrac16$.
Suy ra $abc\le\dfrac16$, $V\ge27$; dấu bằng khi $a=1$, $b=\dfrac12$, $c=\dfrac13$. Giá trị nhỏ nhất bằng $27$.}
\end{ex}
